\section{Number of Islands II}
\label{sec:graph-number-of-islands-2}

\problemstatement{Given an $m \times n$ grid initially entirely water,
and a sequence of positions where land is added \textbf{one cell at a
time}, return the number of islands \textbf{after each addition}.}

\begin{patternbox}
\emph{``Track connected components as elements are added one at a
time''} is exactly what Disjoint Set is built for, in contrast to the
Graph Basics chapter's Number of Islands (a single, static, one-shot
count via traversal): here, the grid \textbf{changes incrementally},
and re-running a full DFS/BFS after every single addition would cost
$O(mn)$ per step. Disjoint Set instead updates the island count in
near-constant time per addition.
\end{patternbox}

\subsection*{Algorithm}

\begin{enumerate}
  \item Maintain a Disjoint Set over all grid cells (indexed as
        $r \times \textit{cols} + c$), and an $\textit{islandCount}$
        starting at $0$.
  \item For each new land position $(r,c)$: if it's already land
        (a duplicate addition), record the current count unchanged and
        continue. Otherwise, mark it land and increment
        $\textit{islandCount}$ by $1$ (a brand new, isolated island).
  \item Check all $4$ neighbors: for each neighbor that is already
        land and in a \textbf{different} Disjoint Set group, union
        them and \textbf{decrement} $\textit{islandCount}$ by $1$
        (two previously separate islands just merged into one).
  \item Record $\textit{islandCount}$ as the answer for this step.
\end{enumerate}

\subsection*{C++ implementation}

\begin{lstlisting}
#include <bits/stdc++.h>
using namespace std;

class DisjointSet {
    vector<int> parent, size;
public:
    DisjointSet(int n) : parent(n), size(n, 1) {
        for (int i = 0; i < n; i++) parent[i] = i;
    }
    int find(int x) {
        if (parent[x] == x) return x;
        return parent[x] = find(parent[x]);
    }
    void unite(int x, int y) {
        int rootX = find(x), rootY = find(y);
        if (rootX == rootY) return;
        if (size[rootX] < size[rootY]) swap(rootX, rootY);
        parent[rootY] = rootX;
        size[rootX] += size[rootY];
    }
};

vector<int> numIslands2(int m, int n, vector<vector<int>>& positions) {
    DisjointSet ds(m * n);
    vector<vector<bool>> isLand(m, vector<bool>(n, false));
    int islandCount = 0;
    vector<int> result;

    int dr[] = {-1, 1, 0, 0};
    int dc[] = {0, 0, -1, 1};

    for (auto& pos : positions) {
        int r = pos[0], c = pos[1];
        if (isLand[r][c]) {
            result.push_back(islandCount);
            continue;
        }

        isLand[r][c] = true;
        islandCount++;

        for (int dir = 0; dir < 4; dir++) {
            int nr = r + dr[dir], nc = c + dc[dir];
            if (nr >= 0 && nr < m && nc >= 0 && nc < n && isLand[nr][nc]) {
                int idA = r * n + c, idB = nr * n + nc;
                if (ds.find(idA) != ds.find(idB)) {
                    ds.unite(idA, idB);
                    islandCount--;
                }
            }
        }
        result.push_back(islandCount);
    }
    return result;
}

int main() {
    vector<vector<int>> positions = {{0,0},{0,1},{1,2},{1,0},{0,2}};
    for (int x : numIslands2(3, 3, positions)) cout << x << " ";
    cout << endl; // 1 1 2 2 1
    return 0;
}
\end{lstlisting}

\subsection*{Dry run}

\dryrun{Take $m=n=3$,
$\textit{positions}=[(0,0),(0,1),(1,2),(1,0),(0,2)]$.}

Add $(0,0)$: new island, count $=1$. No land neighbors. Result: $1$.
Add $(0,1)$: new island, count $=2$; neighbor $(0,0)$ is land, union,
count $=1$. Result: $1$. Add $(1,2)$: new island, count $=2$; no land
neighbors yet. Result: $2$. Add $(1,0)$: new island, count $=3$;
neighbor $(0,0)$ is land (same group as $(0,1)$), union, count $=2$.
Result: $2$. Add $(0,2)$: new island, count $=3$; neighbor $(1,2)$ is
land, union, count $=2$; neighbor $(0,1)$ is land and -- after the
previous union -- in a \textbf{different} group from $(0,2)$'s newly
merged group, union again, count $=1$. Result: $1$. Final sequence:
$1,1,2,2,1$ -- the last addition, $(0,2)$, happens to bridge
\emph{two} previously separate islands at once, dropping the count by
$2$ in a single step.

\begin{complexitybox}
\textbf{Time Complexity.} $O(k \cdot \alpha(mn))$ for $k$ position
updates, each doing a constant number of Disjoint Set operations.
\textbf{Space Complexity.} $O(mn)$ for the Disjoint Set and the
\textit{isLand} grid.
\end{complexitybox}

\begin{takeawaybox}
This problem is the clearest illustration in this chapter of Disjoint
Set's core advantage over re-traversal: when connectivity needs to be
queried or updated \textbf{repeatedly} as a structure changes
incrementally, Disjoint Set answers each update in near-constant time,
where a fresh DFS/BFS per update would be asymptotically far slower.
\end{takeawaybox}
